\documentclass[12pt,a4paper]{article}
\usepackage[utf8]{inputenc}
\usepackage[T2A,T1]{fontenc}
\usepackage[ukrainian]{babel}
\usepackage{amsmath,amssymb}
\usepackage{booktabs}
\usepackage{longtable}
\usepackage{array}
\usepackage{calc}
\usepackage{geometry}
\geometry{margin=2.5cm}

\begin{document}
\section{ОПИС ВИНАХОДУ}\label{ux43eux43fux438ux441-ux432ux438ux43dux430ux445ux43eux434ux443}



\subsection{Назва винаходу}\label{ux43dux430ux437ux432ux430-ux432ux438ux43dux430ux445ux43eux434ux443}



Система та спосіб обробки багатоканальних EEG-даних для обчислення показника обчислювальної ментальної енергії (CME) із застосуванням квантового обчислювального модуля та метаевристичної оптимізації.



\begin{center}\rule{0.5\linewidth}{0.5pt}\end{center}



\subsection{1. Галузь техніки, до якої належить винахід}\label{ux433ux430ux43bux443ux437ux44c-ux442ux435ux445ux43dux456ux43aux438-ux434ux43e-ux44fux43aux43eux457-ux43dux430ux43bux435ux436ux438ux442ux44c-ux432ux438ux43dux430ux445ux456ux434}



\textbf{МПК:} A61B 5/372 (електроенцефалографія); G06N 10/00 (квантові обчислення); G16H 50/20 (медичні ІТ, прогностичне моделювання).



Винахід належить до галузі нейроінформатики, медичних інформаційних технологій та великих даних, зокрема до систем збору та обробки електроенцефалографічних (EEG) даних, методів інтелектуального аналізу часових рядів та архітектур потокової обробки даних.



Винахід також стосується застосування квантових обчислень (квантових процесорів або симуляторів) і варіаційних квантових схем для оцінювання ймовірностей цільових станів, а також метаевристичних алгоритмів оптимізації для налаштування параметрів моделей і ресурсних режимів виконання.



Винахід може бути використано у системах моніторингу психофізіологічного стану людини, нейроінтерфейсах (brain-computer interfaces, BCI), медтех-додатках, системах персоналізованого навчання та продуктивності, а також у дослідницьких платформах для аналізу великих масивів EEG-даних.



\begin{center}\rule{0.5\linewidth}{0.5pt}\end{center}



\subsection{2. Терміни та скорочення}\label{ux442ux435ux440ux43cux456ux43dux438-ux442ux430-ux441ux43aux43eux440ux43eux447ux435ux43dux43dux44f}



У цьому описі застосовано такі терміни та скорочення:



{\def\LTcaptype{none} % do not increment counter

\begin{longtable}[]{@{}

  >{\raggedright\arraybackslash}p{0.4211\linewidth}

  >{\raggedright\arraybackslash}p{0.5789\linewidth}@{}}

\toprule\noalign{}

\begin{minipage}[b]{\linewidth}\raggedright

Термін

\end{minipage} & \begin{minipage}[b]{\linewidth}\raggedright

Визначення

\end{minipage} \\

\midrule\noalign{}

\endhead

\bottomrule\noalign{}

\endlastfoot

EEG-канал (електрод) & Канал реєстрації біоелектричної активності мозку (наприклад, AF7, AF8, TP9, TP10 тощо). \\

\(\Delta\) (вікно аналізу) & Тривалість сегмента сигналу EEG для обчислення ознак. За замовчуванням \(\Delta = 5\) с (конфігурований). Одиниці: с. \\

Діапазони частот & \(\delta\): 1--4 Гц; \(\theta\): 4--8 Гц; \(\alpha\): 8--13 Гц; \(\beta\): 13--30 Гц; \(\gamma\): 30--45 Гц (верхня межа узгоджується з апаратною смугою; \(\gamma\) є опційним --- за наявності апаратної підтримки). \\

\(P_{\text{band}}(ch, t)\) & Інтеграл спектральної густини потужності (PSD) у діапазоні band для електроду \(ch\) у вікні \(t\). Одиниці: мкВ². \\

\(E_{\text{band}}(t)\) & Зведена спектральна потужність EEG-ознак у вікні \(t\) --- зважена сума \(P_{\text{band}}(ch, t)\) по каналах і діапазонах. Одиниці: мкВ². \\

Вн (Вернік; міжнародне: Vernik, Vn) & Скорочене робоче позначення одиниці вимірювання обчислювальної ментальної енергії (CME), введене у цьому описі: \(1\;\text{Вн} \equiv 1\;\text{мкВ}^2 \cdot \text{с}\). Чисельно дорівнює часовому інтегралу квадрату амплітуди EEG-сигналу за вікно спостереження \(\Delta\). Позначення Вн є робочим скороченням, прийнятим у межах цього опису. Типографічне оформлення (прямий шрифт, без крапки) відповідає правилам SI для позначень одиниць {[}23{]}. \\

CME (Computational Mental Energy) & Обчислювальна ментальна енергія --- інтегральний операційно визначений показник, що відображає інтенсивність обробки та залученості користувача за EEG-ознаками у заданому часовому вікні \(\Delta\), обчислений на основі спектральної енергії багатоканальних EEG-ознак та модуляції цієї енергії функцією складності задачі і ймовірності стану потоку, оціненої квантовим модулем. Формально: \(\text{CME}_{\text{rate}}(t) = \kappa \cdot E_{\text{band}}(t) \cdot g(c(t),\,p_{\text{flow}}(t))\) (Вн/с); \(\text{CME}(t) = \text{CME}_{\text{rate}}(t) \cdot \Delta\) (Вн); \(\text{CME}_{\text{session}} = \sum_t \text{CME}(t)\) (Вн). \\

\(\text{CME}_{\text{rate}}(t)\) & Миттєва інтенсивність обчислювальної ментальної енергії у вікні \(t\): \(\text{CME}_{\text{rate}}(t) = \kappa \cdot E_{\text{band}}(t) \cdot g\bigl(c(t),\,p_{\text{flow}}(t)\bigr)\). Одиниці: Вн/с. \\

\(\text{CME}(t)\) & Внесок вікна \(t\) у показник CME: \(\text{CME}(t) = \text{CME}_{\text{rate}}(t) \cdot \Delta\). Одиниці: Вн. \\

\(\text{CME}_{\text{session}}\) & Сумарний показник CME за всю сесію: \(\sum_t \text{CME}(t)\). Одиниці: Вн. \\

\(\text{CME}_{\text{index}}(t)\) & Безрозмірна шкала відображення CME \([0,\,100]\), отримана нормалізацією \(\text{CME}_{\text{rate}}(t)\) або \(\text{CME}(t)\) за калібрувальними даними. Не є базовою одиницею. \\

\(\text{CME}_J(t)\) & Факультативне наближене відображення CME у систему одиниць SI: \(\text{CME}_J(t) = 10^{-12} / Z_e \cdot \text{CME}(t)\), де \(Z_e\) --- імпеданс електрода (Ом). Одиниці: Дж. Є апроксимацією, що залежить від електричної моделі та вимірюваного \(Z_e\); не є основним результатом системи. \\

\(\mathbf{x}_t\) & Повний вхідний вектор ознак у вікні \(t\) (\(N_{ch} \times N_{band} + 2\) ознак; для 4 електродів × 5 діапазонів: 22 ознаки; за відсутності \(\gamma\): 4 × 4 + 2 = 18): спектральні потужності по кожному електроду, похідні індекси, \(c(t)\), \(q(t)\). \\

\(\mathbf{x}_t^{(q)}\) & Зведений вектор ознак для квантового модуля (\(\in \mathbb{R}^8\)): середні потужності \(\overline{P}_{\delta/\theta/\alpha/\beta/\gamma}\), фронтальна \(\alpha\)-асиметрія, індекс залученості \(\beta/\theta\), складність задачі \(c(t)\). \\

QPU & Квантовий обчислювальний модуль (реальний квантовий процесор) або симулятор квантових обчислень. \\

shots (\(S\)) & Кількість вимірювань квантової схеми для оцінки ймовірностей результатів. За замовчуванням \(S = 1024\). Одиниці: шт. \\

Ansatz, \(D\) & Варіаційна структура квантової схеми; \(D\) --- глибина (кількість шарів). \\

Стан потоку (flow state) & Психологічний стан повного занурення у діяльність з оптимальним балансом між складністю задачі та навичками виконавця {[}1{]}. У цьому описі стан потоку визначається операційно через EEG-кореляти {[}4, 7{]} і є цільовим класом бінарної класифікації квантовою моделлю. Валідація виконується за визнаними психометричними шкалами {[}24, 25{]}. У формулах позначається як ``flow''. \\

\(p_{\text{flow}}(t)\) & Ймовірність стану потоку (flow) у вікні \(t\), оцінена квантовим модулем як вихід навченого класифікатора (п. 6.7.3). Позначення \(p_{\text{flow}}(t)\) є базовим у формулах; альтернативне позначення \(p_\text{focus}(t)\) може використовуватися у програмних інтерфейсах. Одиниці: безрозмірна, \(\in [0,\,1]\). \\

Метаевристика & Алгоритм оптимізації (GA/PSO/ACO/SA тощо) для підбору параметрів моделі та/або режимів виконання. \\

\(c(t)\) & Нормований індекс складності/навантаження у вікні \(t\). Одиниці: безрозмірний, \(\in [0,\,1]\). \\

\(q(t)\) & Індикатор якості сигналу у вікні \(t\); може враховувати імпеданс контакту електроду \(Z_{ch}(t)\) для зважування артефактів. Одиниці: безрозмірний, \(\in [0,\,1]\). \\

GA & Генетичний алгоритм (Genetic Algorithm). \\

PSO & Оптимізація роєм частинок (Particle Swarm Optimization). \\

ACO & Мурашиний алгоритм (Ant Colony Optimization). \\

SA & Імітація відпалу (Simulated Annealing). \\

VQC & Варіаційний квантовий класифікатор (Variational Quantum Classifier). \\

QSVC & Квантова машина опорних векторів (Quantum Support Vector Classifier). \\

fMRI & Функціональна магнітно-резонансна томографія (не є частиною основного EEG-пайплайну; згадується лише як потенційне зовнішнє джерело для валідації). \\

\(T_{\text{qpu}}\) & Латентність виконання квантової схеми на QPU/симуляторі; може бути виміряною (реальний QPU) або змодельованою (симулятор) величиною. Одиниці: мс. \\

Пристрій стримінгу & Програмно-апаратний компонент (наприклад, Raspberry Pi), що приймає дані EEG від пристрою реєстрації та передає їх на сервер. \\

\end{longtable}

}



\begin{center}\rule{0.5\linewidth}{0.5pt}\end{center}



\subsection{3. Рівень техніки та найближчий аналог}\label{ux440ux456ux432ux435ux43dux44c-ux442ux435ux445ux43dux456ux43aux438-ux442ux430-ux43dux430ux439ux431ux43bux438ux436ux447ux438ux439-ux430ux43dux430ux43bux43eux433}



\subsubsection{3.1. Відомі рішення}\label{ux432ux456ux434ux43eux43cux456-ux440ux456ux448ux435ux43dux43dux44f}



Оцінювання когнітивного навантаження та стану потоку (flow state) --- стану повного занурення у діяльність з оптимальним балансом між складністю задачі та навичками виконавця {[}1{]} --- є одним із пріоритетних завдань прикладної нейроінформатики. Канонічним прикладом кількісного оцінювання когнітивної залученості є індекс \(\beta/(\alpha+\theta)\), запропонований як показник операторської залученості {[}2{]} та підтверджений у задачах адаптивної автоматизації {[}3{]}. Відношення \(\alpha/\theta\) та \(\theta/\alpha\) також використовуються для дискримінації рівнів когнітивного навантаження {[}6{]}. Нейрофізіологічні кореляти стану потоку досліджуються за EEG: зокрема, підвищення фронтальної \(\theta\)-активності та помірна \(\alpha\)-активність у лобово-центральних ділянках асоціюються зі станом потоку при когнітивних завданнях {[}4{]}; об'єктивне оцінювання потоку за портативними EEG-пристроями запропоновано як крок від суб'єктивних шкал до замкнутого контуру зворотного зв'язку {[}7{]}; також досліджено одноканальні префронтальні EEG-кореляти потоку {[}8{]} та зв'язок EEG-індексів залученості із самозвітним рівнем потоку у когнітивних іграх {[}9{]}.



Відомі системи оцінювання когнітивних, емоційних або психофізіологічних станів за EEG {[}16, 17{]} використовують цифрову обробку сигналів та класичні методи машинного навчання для класифікації/регресії, зокрема на основі спектральних ознак у діапазонах \(\delta/\theta/\alpha/\beta/\gamma\). Такі рішення зазвичай реалізують ланцюжок:



\begin{quote}

зчитування EEG - сегментація на епохи - обчислення ознак - класичний класифікатор - індикатор стану.

\end{quote}



Окремо відомі підходи квантового машинного навчання, де варіаційні квантові класифікатори {[}5, 14, 15{]} або квантові SVM застосовують до задач класифікації EEG {[}10, 11{]}. Також відомі метаевристичні оптимізатори для налаштування параметрів моделей та керування ресурсами квантових обчислень {[}12, 13{]}.



\subsubsection{3.2. Недоліки відомих рішень}\label{ux43dux435ux434ux43eux43bux456ux43aux438-ux432ux456ux434ux43eux43cux438ux445-ux440ux456ux448ux435ux43dux44c}



Однак відомі рішення, як правило, не забезпечують єдиного наскрізного процесу, що поєднує: (i) стандартизовану сегментацію EEG у фіксовані часові вікна; (ii) обчислення багатоканальних спектральних ознак по кожному електроду в усіх частотних діапазонах; (iii) квантову оцінку ймовірності цільового стану за допомогою параметризованої варіаційної схеми; (iv) обчислення інтегрального показника CME з урахуванням енергетичних, когнітивних та квантових компонент; (v) одночасне керування ресурсами QPU (shots, глибина, латентність) через метаевристичну оптимізацію.



\subsubsection{3.3. Найближчий аналог}\label{ux43dux430ux439ux431ux43bux438ux436ux447ux438ux439-ux430ux43dux430ux43bux43eux433}



Серед патентних публікацій відомі рішення для детектування психічних станів за EEG {[}16{]}, самокалібрувального нейрозворотного зв'язку {[}17{]}, досягнення стану потоку за допомогою біозворотного зв'язку від портативних сенсорів {[}18{]} та визначення і кореляції стану потоку з контекстом діяльності {[}19{]}. У галузі керування ресурсами QPU відомі: оптимізація квантових завдань у черзі {[}20{]}, спільне планування квантових задач з урахуванням обмежень {[}21{]} та витіснення квантових програм за кількістю пострілів (shots) {[}22{]}.



Найближчим аналогом обрано систему детектування ментальних станів за EEG-даними US7865235B2 {[}16{]}, яка реалізує ланцюжок ``зчитування EEG --- спектральні ознаки --- класичний класифікатор --- індикатор стану'' і є найбільш повним рішенням серед відомих за сукупністю ознак. Ця система: (i) не містить квантового обчислювального модуля; (ii) не містить метаевристичної оптимізації параметрів квантової частини і ресурсних режимів виконання; (iii) не забезпечує обчислення інтегрального показника CME як функції одночасно спектральної енергії, складності задачі та квантово-оціненої ймовірності стану потоку.



\begin{center}\rule{0.5\linewidth}{0.5pt}\end{center}



\subsection{4. Суть винаходу}\label{ux441ux443ux442ux44c-ux432ux438ux43dux430ux445ux43eux434ux443}



\subsubsection{4.1. Задача винаходу}\label{ux437ux430ux434ux430ux447ux430-ux432ux438ux43dux430ux445ux43eux434ux443}



Задачею винаходу є підвищення точності та відтворюваності оцінювання стану користувача за багатоканальними EEG-даними у потоковому режимі, а також забезпечення керованості ресурсів квантового обчислювального модуля при обробці великих обсягів даних.



\subsubsection{4.2. Технічний результат}\label{ux442ux435ux445ux43dux456ux447ux43dux438ux439-ux440ux435ux437ux443ux43bux44cux442ux430ux442}



Технічний результат полягає у формуванні стандартизованого вимірюваного показника обчислювальної ментальної енергії CME --- інтегральної операційно визначеної величини, що відображає інтенсивність обробки та залученості користувача за EEG-ознаками у часовому вікні \(\Delta\) на основі спектральної енергії багатоканальних EEG-ознак, модульованої функцією складності задачі та ймовірності стану потоку. Показник CME вимірюється у одиницях Вн (\(1\;\text{Вн} \equiv 1\;\text{мкВ}^2 \cdot \text{с}\), робоче позначення, прийняте у цьому описі) та обчислюється із залученням квантової моделі для оцінювання ймовірності цільового стану \(p_{\text{flow}}(t)\) та метаевристичної оптимізації параметрів квантової частини з урахуванням якості оцінювання та обчислювальної вартості (shots, глибина, латентність).



Метаевристична оптимізація ресурсних параметрів забезпечує керованість обчислювальних витрат та баланс між якістю оцінювання і вартістю звернень до QPU. Інтеграція запропонованих технологій підтримує горизонтальне масштабування та обробку великих архівів EEG-даних через відповідну архітектуру розгортання (п. 6.11).



\subsubsection{4.3. Ключові відмінні ознаки}\label{ux43aux43bux44eux447ux43eux432ux456-ux432ux456ux434ux43cux456ux43dux43dux456-ux43eux437ux43dux430ux43aux438}



Ключові відмінні ознаки запропонованого рішення:



\begin{enumerate}

\def\labelenumi{\arabic{enumi}.}

\item

  Збереження повної багатоканальної інформації на етапі побудови первинних ознак (потужності у діапазонах \(\delta/\theta/\alpha/\beta/\gamma\) по кожному електроду); допускається подальша агрегація/компресія для формування зведеного вектора квантового модуля.

\item

  Квантова оцінка стану через формування 8-ознакового вхідного вектора \(\mathbf{x}_t^{(q)}\) та подання його у 4-кубітну варіаційну квантову схему з архітектурою повторного завантаження даних (data re-uploading), ZZ-взаємодіями ознак та кільцевим заплутуванням, з отриманням ймовірності \(p_{\text{flow}}(t)\) на QPU або симуляторі.

\item

  Обчислення \(\text{CME}(t)\) з визначеною одиницею Вн (\(1\;\text{Вн} \equiv 1\;\text{мкВ}^2 \cdot \text{с}\)) як функції зведеної спектральної потужності \(E_{\text{band}}(t)\), індексу складності \(c(t)\) та ймовірності стану потоку \(p_{\text{flow}}(t)\) у вікні \(\Delta\), із факультативним відображенням у SI через імпеданс електрода \(Z_e\).

\item

  Метаевристична оптимізація для підбору параметрів квантової схеми \(\Theta\) та ресурсних параметрів \((S, D)\) за цільовою функцією, що включає якість класифікації і ресурсну вартість, із підтримкою алгоритмів GA, PSO, ACO, SA.

\item

  Підтримка монолітного, мікросервісного та брокеризованого режимів розгортання без зміни математичної суті процесу.

\item

  Спосіб взаємодії пристрою реєстрації нейросигналів (EEG), пристрою стримінгу та серверного модуля обчислення CME --- пристрій реєстрації передає дані по Bluetooth на пристрій стримінгу (наприклад, Raspberry Pi), який буферизує, сегментує та передає на сервер через мережевий протокол; сервер обробляє потік, викликає квантовий модуль та обчислює CME.

\item

  Система виконує адаптивний вибір режиму інференсу для кожного вікна на основі якості сигналу \(q(t)\), виміряної або оціненої латентності QPU \(T_{\text{qpu}}(t)\) та ресурсного бюджету. \(p_{\text{flow}}(t)\) може бути отримано від QPU, симулятора або наближеного класичного сурогату з подальшою корекцією при наступному зверненні до QPU.

\end{enumerate}



\begin{center}\rule{0.5\linewidth}{0.5pt}\end{center}



\subsection{5. Опис креслень}\label{ux43eux43fux438ux441-ux43aux440ux435ux441ux43bux435ux43dux44c}



\textbf{Фіг. 1} --- Загальна структурна схема системи (позиції 100--210): пристрій реєстрації нейросигналів EEG (110), пристрій стримінгу (115), модуль приймання даних на сервері (120), модуль сегментації (130), модуль попередньої обробки (140), модуль виділення спектральних ознак (150), модуль формування вектора ознак (160), квантовий обчислювальний модуль (170), модуль обчислення CME (180), модуль метаевристичної оптимізації (190), модуль класичного нейромережевого аналізу (195), сховище/журнал (200), модуль асинхронного збереження даних (205), інтерфейс інтеграції (210).



\textbf{Фіг. 2} --- Потік даних для онлайн-обчислення CME у вікнах \(\Delta\). Показано послідовність операцій від надходження EEG-даних до видачі значення \(\text{CME}(t)\) із зазначенням взаємодії з квантовим модулем.



\textbf{Фіг. 3} --- Контур метаевристичної оптимізації параметрів квантової моделі та ресурсних режимів (shots/глибина). Показано ітераційний цикл: ініціалізація популяції - оцінка кандидатів через QPU - обчислення фітнес-функції - еволюція популяції - оновлення найкращих параметрів.



\textbf{Фіг. 4} --- Приклад структури спектральних ознак по електродах та діапазонах частот у межах одного вікна. Ілюстрація формування вектора \(\mathbf{x}_t\) із потужностей \(\delta/\theta/\alpha/\beta/\gamma\) для кожного каналу.



\begin{center}\rule{0.5\linewidth}{0.5pt}\end{center}



\subsection{6. Опис винаходу}\label{ux43eux43fux438ux441-ux432ux438ux43dux430ux445ux43eux434ux443-1}



\subsubsection{6.1. Склад та статична структура системи}\label{ux441ux43aux43bux430ux434-ux442ux430-ux441ux442ux430ux442ux438ux447ux43dux430-ux441ux442ux440ux443ux43aux442ux443ux440ux430-ux441ux438ux441ux442ux435ux43cux438}



На Фіг. 1 зображено систему (100), яка містить такі модулі:



\textbf{(110) Пристрій реєстрації нейросигналів (EEG)} --- апаратний засіб реєстрації біоелектричної активності мозку (наприклад, Muse Athena, Emotiv, OpenBCI або аналогічний пристрій із щонайменше 4 електродами EEG: TP9, AF7, AF8, TP10); з'єднання з пристроєм стримінгу через Bluetooth або інший бездротовий протокол.



\textbf{(115) Пристрій стримінгу} --- програмно-апаратний компонент (наприклад, одноплатний комп'ютер типу Raspberry Pi), що приймає сирі дані EEG від пристрою реєстрації (110) через Bluetooth, буферизує їх, виконує попередню сегментацію та передає на сервер через мережевий протокол (OSC, WebSocket, SignalR, HTTP або інший).



\textbf{(120) Модуль приймання даних} --- програмний компонент на сервері, що отримує сирі або попередньо оброблені потоки EEG від пристрою стримінгу (115) через двонаправлений потоковий протокол (наприклад, WebSocket, SignalR) або REST API; зберігає часові мітки та метадані.



\textbf{(130) Модуль сегментації} --- виконує поділ безперервного потоку EEG-даних на дискретні часові вікна тривалості \(\Delta\). Для \(k\)-го вікна:



\[W_k = \bigl[k \cdot \Delta,\;\;(k+1) \cdot \Delta\bigr)\]



Кожне вікно представлене структурою даних, що містить часові мітки початку та кінця вікна, а тривалість вікна обчислюється як їх різниця.



\textbf{(140) Модуль попередньої обробки} --- виконує фільтрацію та нормалізацію сигналу, включно з пригніченням мережевих завад (50 Гц), видаленням дрейфу, контролем якості контакту електродів та маркуванням артефактів. Нормалізація ознак виконується відповідним програмним модулем, який масштабує кожну ознаку до діапазону \([-1,\,1]\) методом мін-макс перетворення з калібрувальними статистиками:



\[f_{\text{norm}}[i] = 2 \cdot \frac{f[i] - f_{\min}^{(i)}}{f_{\max}^{(i)} - f_{\min}^{(i)}} - 1\]



де \(f_{\min}^{(i)}\), \(f_{\max}^{(i)}\) --- мінімальне та максимальне значення ознаки \(i\), визначені під час калібрувального періоду (перші \(N_{\text{cal}}\) вікон сесії або з персонального профілю). За відсутності калібрувальних даних може використовуватись z-нормалізація з подальшим обмеженням \(\tanh\). Діапазон \([-1,\,1]\) забезпечує повне покриття кутів обертання \([0,\,2\pi]\) при кодуванні у квантовій схемі (див. п. 6.7). Параметри фільтрації узгоджуються з апаратними характеристиками EEG-пристрою.



\textbf{(150) Модуль виділення спектральних ознак} --- для кожного електроду \(ch\) та вікна \(t\) обчислює спектральну потужність у стандартних діапазонах: \(P_\delta(ch,t)\), \(P_\theta(ch,t)\), \(P_\alpha(ch,t)\), \(P_\beta(ch,t)\) та опційно \(P_\gamma(ch,t)\). Обчислення виконується за допомогою FFT/PSD або іншого узгодженого методу. \(\gamma\)-діапазон (30--45 Гц) включається за наявності апаратної підтримки відповідної смуги пропускання. У реалізації спектральні потужності зберігаються як числові поля відповідної структури даних. Одиниці вимірювання: мкВ² (інтеграл PSD по діапазону).



\textbf{(160) Модуль формування вектора ознак} --- формує вектор \(\mathbf{x}_t\), що містить спектральні потужності у діапазонах \(\delta/\theta/\alpha/\beta/\gamma\) по кожному електроду, а також похідні індекси та метадані задачі:



\[\mathbf{x}_t = \Bigl[\,P_\delta(ch,t),\;P_\theta(ch,t),\;P_\alpha(ch,t),\;P_\beta(ch,t),\;P_\gamma(ch,t)\;\Big|\;\forall\,ch \in \mathcal{CH}\;;\;c(t)\;;\;q(t)\,\Bigr]\]



де \(c(t)\) --- нормований індекс складності/навантаження у вікні \(t\) (за замовчуванням = 0.5), \(q(t)\) --- індикатори якості сигналу (включаючи, за наявності, дані про імпеданс контакту електроду \(Z_{ch}(t)\) для зниження впливу артефактів).



\textbf{(170) Квантовий обчислювальний модуль} --- реалізує 4-кубітну варіаційну квантову схему з параметрами \(\Theta\) та глибиною \(D\). Деталі квантової схеми описані у п. 6.7.



\textbf{(180) Модуль обчислення CME} --- обчислює показник \(\text{CME}(t)\) (одиниці: Вн) за формулою, описаною у п. 6.8.



\textbf{(190) Модуль метаевристичної оптимізації} --- фоновий програмний сервіс, що виконує підбір параметрів \(\Theta\) квантової схеми та ресурсних параметрів \((S, D)\) за цільовою функцією із підтримкою алгоритмів GA, PSO, ACO, SA.



\textbf{(195) Модуль класичного нейромережевого аналізу} --- окремий програмний сервіс з навченою класичною нейронною мережею для оцінювання стану потоку за 22-ознаковим вектором (20 спектральних потужностей по 4 електродах у 5 діапазонах + складність задачі + якість сигналу). Модуль повертає ймовірність стану потоку \(p_{\text{flow}}^{\text{NN}}(t) \in [0,\,1]\) та бінарну мітку \(\text{FlowLabel}\). Використовується для масового маркування набору даних, порівняльного аналізу з квантовою моделлю та гібридного інференсу.



\textbf{(200) Сховище/журнал} --- реляційна база даних (наприклад, SQL Server або аналогічна) для зберігання сесій, журналів запитів, результатів CME по вікнах, завдань навчання, журналів звернень до QPU, а також набору даних EEG-ознак по вікнах із мітками стану потоку та маркерів дій користувача.



\textbf{(205) Модуль асинхронного збереження даних} --- фоновий програмний сервіс, що приймає потокові EEG-вікна через обмежену чергу і асинхронно записує їх у сховище пакетами. Це забезпечує неблокуючу обробку: дані зберігаються незалежно від результатів інференсу або обчислення CME. При записі автоматично призначається евристична мітка стану потоку на основі відношення фронтальної \(\alpha/\theta\) потужності. Модуль автоматично прив'язує вікна до маркерів дій за часовою належністю.



\textbf{(210) Інтерфейс інтеграції} --- REST API та веб-інтерфейс для моніторингу, подання запитів, управління навчанням, візуалізації результатів, експорту даних, керування сесіями збору даних (старт/стоп), запуску класичного аналізу та візуалізації розподілу міток стану потоку.



У статичному стані зазначені модулі можуть бути реалізовані як програмні компоненти одного застосунку або як окремі сервіси/воркери у розподіленій архітектурі (див. п. 6.10).



\subsubsection{6.2. Вхідні дані та сегментація}\label{ux432ux445ux456ux434ux43dux456-ux434ux430ux43dux456-ux442ux430-ux441ux435ux433ux43cux435ux43dux442ux430ux446ux456ux44f}



Система приймає потоки багатоканальних EEG-даних із часовими мітками та метаданими. Основні вхідні дані можуть надходити у реальному часі від пристрою реєстрації EEG через пристрій стримінгу (наприклад, Raspberry Pi, підключений до Muse Athena по Bluetooth) або у вигляді CSV-файлів з попередньо зареєстрованими записами чи експорту з системи.



Формат запису для одного вікна включає такі складові:



{\def\LTcaptype{none} % do not increment counter

\begin{longtable}[]{@{}

  >{\raggedright\arraybackslash}p{0.6250\linewidth}

  >{\raggedright\arraybackslash}p{0.3750\linewidth}@{}}

\toprule\noalign{}

\begin{minipage}[b]{\linewidth}\raggedright

Складова

\end{minipage} & \begin{minipage}[b]{\linewidth}\raggedright

Опис

\end{minipage} \\

\midrule\noalign{}

\endhead

\bottomrule\noalign{}

\endlastfoot

Часова мітка & timestamp \\

Ідентифікатор сесії & session\_id \\

Ідентифікатор вікна & window\_id \\

Спектральні потужності & по діапазонах \(\delta/\theta/\alpha/\beta/\gamma\) та електродах \\

Похідні ознаки & frontal\_asym, parietal\_asym, engagement \\

Складність задачі & task\_difficulty \\

Мітка стану & Flow / No\_Flow (для навчальних даних) \\

\end{longtable}

}



Сегментація здійснюється у вікна тривалості \(\Delta\). Для \(k\)-го вікна: \(W_k = [k \cdot \Delta,\;(k+1) \cdot \Delta)\). У програмній реалізації тривалість вікна визначається різницею часових міток кінця та початку відповідної структури даних, із значенням за замовчуванням 5 секунд.



\subsubsection{6.3. Попередня обробка}\label{ux43fux43eux43fux435ux440ux435ux434ux43dux44f-ux43eux431ux440ux43eux431ux43aux430}



Модуль попередньої обробки (140) виконує фільтрацію та нормалізацію сигналу, пригнічення мережевих завад (50 Гц), видалення дрейфу, контроль якості контакту електродів та маркування артефактів.



У програмній реалізації попередня обробка включає нормалізацію вектора ознак до діапазону \([-1,\,1]\):



\[f_{\text{norm}}[i] = 2 \cdot \frac{f[i] - f_{\min}^{(i)}}{f_{\max}^{(i)} - f_{\min}^{(i)}} - 1\]



де \(f_{\min}^{(i)}\), \(f_{\max}^{(i)}\) --- мінімальне та максимальне значення ознаки \(i\) з калібрувального набору або персонального профілю. За відсутності калібрувальних даних застосовується z-нормалізація з подальшим обмеженням \(\tanh\):



\[f_{\text{norm}}[i] = \tanh\!\left(\frac{f[i] - \mu_i}{\sigma_i}\right)\]



Діапазон \([-1,\,1]\) забезпечує кути обертання у \([0,\,2\pi]\) при кодуванні \((f_{\text{norm}} + 1) \cdot \pi\), що покриває повну сферу Блоха (див. п. 6.7).



\subsubsection{6.4. Частотні ознаки по кожному електроду}\label{ux447ux430ux441ux442ux43eux442ux43dux456-ux43eux437ux43dux430ux43aux438-ux43fux43e-ux43aux43eux436ux43dux43eux43cux443-ux435ux43bux435ux43aux442ux440ux43eux434ux443}



Для кожного електроду \(ch\) та вікна \(t\) обчислюється спектральна потужність у стандартних діапазонах:



\[P_\delta(ch, t),\quad P_\theta(ch, t),\quad P_\alpha(ch, t),\quad P_\beta(ch, t),\quad P_\gamma(ch, t)\]



Обчислення PSD виконується методом Уелча (Welch) або еквівалентним спектральним оцінювачем. Базова верхня межа \(\gamma\)-діапазону становить 45 Гц; за наявності апаратної підтримки вона може бути розширена до 50 Гц.



У реалізації спектральні потужності по кожному каналу зберігаються як числові поля відповідної структури даних:



{\def\LTcaptype{none} % do not increment counter

\begin{longtable}[]{@{}

  >{\raggedright\arraybackslash}p{0.1875\linewidth}

  >{\raggedright\arraybackslash}p{0.3125\linewidth}

  >{\raggedright\arraybackslash}p{0.1875\linewidth}

  >{\raggedright\arraybackslash}p{0.3125\linewidth}@{}}

\toprule\noalign{}

\begin{minipage}[b]{\linewidth}\raggedright

Поле

\end{minipage} & \begin{minipage}[b]{\linewidth}\raggedright

Діапазон

\end{minipage} & \begin{minipage}[b]{\linewidth}\raggedright

Межі

\end{minipage} & \begin{minipage}[b]{\linewidth}\raggedright

Примітка

\end{minipage} \\

\midrule\noalign{}

\endhead

\bottomrule\noalign{}

\endlastfoot

\(P_\delta\) & \(\delta\) & 1--4 Гц & \\

\(P_\theta\) & \(\theta\) & 4--8 Гц & \\

\(P_\alpha\) & \(\alpha\) & 8--13 Гц & \\

\(P_\beta\) & \(\beta\) & 13--30 Гц & \\

\(P_\gamma\) & \(\gamma\) & 30--45 Гц & опційне; за наявності апаратної підтримки \\

\end{longtable}

}



Для багатоканального запису від пристрою стримінгу ознаки подаються окремо для кожного електроду (наприклад, Delta\_TP9, Theta\_AF7, Alpha\_AF8, Beta\_TP10 тощо).



\subsubsection{\texorpdfstring{6.5. Зведена спектральна потужність \(E_{\text{band}}(t)\)}{6.5. Зведена спектральна потужність E\_\{\textbackslash text\{band\}\}(t)}}\label{ux437ux432ux435ux434ux435ux43dux430-ux441ux43fux435ux43aux442ux440ux430ux43bux44cux43dux430-ux43fux43eux442ux443ux436ux43dux456ux441ux442ux44c-e_textbandt}



Зведена спектральна потужність \(E_{\text{band}}(t)\) визначається як зважена сума спектральних потужностей по каналах та діапазонах:



\[E_{\text{band}}(t) = \sum_{ch \in \mathcal{CH}} \left[\,w_\delta \cdot P_\delta(ch,t) + w_\theta \cdot P_\theta(ch,t) + w_\alpha \cdot P_\alpha(ch,t) + w_\beta \cdot P_\beta(ch,t) + w_\gamma \cdot P_\gamma(ch,t)\,\right]\]



Оскільки \(P_{\text{band}}(ch, t)\) вимірюється у мкВ², а ваги \(w_\delta,\;w_\theta,\;w_\alpha,\;w_\beta,\;w_\gamma\) є безрозмірними, величина \(E_{\text{band}}(t)\) має одиниці мкВ².



Коефіцієнти задаються попередньо або підбираються під час калібрування. Складова \(w_\gamma \cdot P_\gamma\) є опційною: вона включається за наявності апаратної підтримки \(\gamma\)-діапазону (смуга пропускання не менше 45 Гц). За відсутності \(\gamma\)-даних встановлюється \(w_\gamma = 0\).



У програмній реалізації базова формула для одного каналу обчислюється відповідним програмним модулем:



\[E_{\text{band}}(t) = w_\delta \cdot P_\delta + w_\theta \cdot P_\theta + w_\alpha \cdot P_\alpha + w_\beta \cdot P_\beta + w_\gamma \cdot P_\gamma\]



Значення ваг за замовчуванням:



{\def\LTcaptype{none} % do not increment counter

\begin{longtable}[]{@{}

  >{\raggedright\arraybackslash}p{0.2308\linewidth}

  >{\raggedright\arraybackslash}p{0.3846\linewidth}

  >{\raggedright\arraybackslash}p{0.3846\linewidth}@{}}

\toprule\noalign{}

\begin{minipage}[b]{\linewidth}\raggedright

Вага

\end{minipage} & \begin{minipage}[b]{\linewidth}\raggedright

Значення

\end{minipage} & \begin{minipage}[b]{\linewidth}\raggedright

Примітка

\end{minipage} \\

\midrule\noalign{}

\endhead

\bottomrule\noalign{}

\endlastfoot

\(w_\delta\) & 0.5 & \\

\(w_\theta\) & 1.0 & \\

\(w_\alpha\) & 1.0 & \\

\(w_\beta\) & 0.3 & \\

\(w_\gamma\) & 0.0 & опційно; встановлюється \textgreater{} 0 при наявності \(\gamma\)-даних \\

\end{longtable}

}



Ваги можуть бути змінені через файл конфігурації.



\subsubsection{\texorpdfstring{6.6. Формування вхідного вектора ознак \(\mathbf{x}_t\)}{6.6. Формування вхідного вектора ознак \textbackslash mathbf\{x\}\_t}}\label{ux444ux43eux440ux43cux443ux432ux430ux43dux43dux44f-ux432ux445ux456ux434ux43dux43eux433ux43e-ux432ux435ux43aux442ux43eux440ux430-ux43eux437ux43dux430ux43a-mathbfx_t}



Модуль (160) формує вектор \(\mathbf{x}_t\), який містить спектральні потужності у діапазонах \(\delta/\theta/\alpha/\beta/\gamma\) по кожному електроду, похідні індекси (frontal\_asym, parietal\_asym, engagement), індекс складності задачі \(c(t)\) та індикатори якості сигналу \(q(t)\).



У загальному вигляді:



\[\mathbf{x}_t = \Bigl[\,P_\delta(ch,t),\;P_\theta(ch,t),\;P_\alpha(ch,t),\;P_\beta(ch,t),\;P_\gamma(ch,t)\;\Big|\;\forall\,ch \in \mathcal{CH}\;;\;c(t)\;;\;q(t)\,\Bigr]\]



Для 4-кубітної квантової схеми вхідний вектор зводиться до 8 ознак:



\[\mathbf{x}_t^{(q)} = \Bigl[\,\overline{P}_\delta,\;\overline{P}_\theta,\;\overline{P}_\alpha,\;\overline{P}_\beta,\;\overline{P}_\gamma,\;\text{FrontalAsym}_\alpha,\;\text{Engagement}_{\beta/\theta},\;c(t)\,\Bigr] \in \mathbb{R}^8\]



де \(\overline{P}_{\text{band}}\) --- середнє спектральної потужності по всіх каналах, \(\text{FrontalAsym}_\alpha = P_\alpha(\text{AF8}) - P_\alpha(\text{AF7})\) --- фронтальна \(\alpha\)-асиметрія, \(\text{Engagement}_{\beta/\theta} = \overline{P}_\beta / \overline{P}_\theta\) --- індекс залученості, \(c(t)\) --- складність задачі. Якщо \(\gamma\)-діапазон недоступний (див. п. 6.5), \(\overline{P}_\gamma = 0\), при цьому розмірність вектора \(\mathbb{R}^8\) зберігається.



Вхідний вектор передається до квантового модуля (170) через мережевий протокол обміну даними (приклад формату запиту наведено у п. 7.3).



Додатково, при наявності навчених параметрів моделі, вектор \texttt{trainedParams} (\(8 \times (L + 1) = 24\) чисел для \(L = 2\) шарів повторного завантаження) передається разом з ознаками для використання у варіаційних шарах квантової схеми.



\subsubsection{6.7. Квантова оцінка ймовірності цільового стану}\label{ux43aux432ux430ux43dux442ux43eux432ux430-ux43eux446ux456ux43dux43aux430-ux439ux43cux43eux432ux456ux440ux43dux43eux441ux442ux456-ux446ux456ux43bux44cux43eux432ux43eux433ux43e-ux441ux442ux430ux43dux443}



Квантовий модуль (170) реалізує 4-кубітну варіаційну квантову схему з архітектурою повторного завантаження даних (data re-uploading) та ZZ-взаємодіями ознак. Модуль реалізований як окремий програмний сервіс із використанням квантового симулятора (наприклад, Qiskit Aer). За потреби, модуль може бути підключений до реального квантового процесора.



\paragraph{6.7.1. Архітектура повторного завантаження даних}\label{ux430ux440ux445ux456ux442ux435ux43aux442ux443ux440ux430-ux43fux43eux432ux442ux43eux440ux43dux43eux433ux43e-ux437ux430ux432ux430ux43dux442ux430ux436ux435ux43dux43dux44f-ux434ux430ux43dux438ux445}



На відміну від одноразового кодування ознак, архітектура повторного завантаження (data re-uploading) {[}5{]} повторно кодує класичні дані у квантовий стан \(L + 1\) разів, чергуючи кодування з навченими варіаційними перетвореннями. Це підвищує виразність схеми (expressivity) {[}14{]}, оскільки число частот у Фур'є-спектрі квантової моделі зростає пропорційно до кількості повторних завантажень. Для \(L = 2\) шарів повторного завантаження та \(N_q = 4\) кубітів загальна кількість навчених параметрів становить:



\[|\Theta| = 2 \cdot N_q \cdot (L + 1) = 2 \times 4 \times 3 = 24\]



\paragraph{6.7.2. Структура одного шару}\label{ux441ux442ux440ux443ux43aux442ux443ux440ux430-ux43eux434ux43dux43eux433ux43e-ux448ux430ux440ux443}



Кожен шар \(l \in \{0, 1, \ldots, L\}\) містить чотири підшари:



Підшар A (двоканальне кутове кодування). 8 ознак кодуються у квантовий стан через обертання \(R_y\) та \(R_z\) --- по дві осі обертання на кубіт, що дозволяє закодувати удвічі більше ознак порівняно з одноканальним кодуванням:



\[\forall\,i \in \{0,1,2,3\}:\quad R_y\bigl((x_i + 1) \cdot \pi\bigr)\,|q_i\rangle, \quad R_z\bigl((x_{i+4} + 1) \cdot \pi\bigr)\,|q_i\rangle\]



Припускається, що ознаки нормалізовані до діапазону \([-1,\,1]\), що забезпечує кути обертання у діапазоні \([0,\,2\pi]\).



Підшар B (ZZ-взаємодія ознак). Для захоплення нелінійних взаємодій між ознаками застосовуються ZZ-гейти (реалізовані через CNOT-RZ-CNOT декомпозицію) для кільцевих пар кубітів:



\[\forall\,i \in \{0,1,2,3\},\;j = (i+1) \bmod 4:\quad RZZ\bigl((x_i - x_j)^2 \cdot \pi,\;q_i,\;q_j\bigr)\]



де \(RZZ(\phi, q_i, q_j) = \text{CNOT}(q_i, q_j) \cdot R_z(\phi, q_j) \cdot \text{CNOT}(q_i, q_j)\).



Ця конструкція аналогічна ZZFeatureMap {[}15{]} і забезпечує кодування попарних кореляцій між ознаками у квантовий стан.



Підшар C (кільцеве заплутування). Застосовуються CNOT-гейти у кільцевій топології:



\[\forall\,i \in \{0,1,2,3\}:\quad \text{CNOT}(q_i,\,q_{(i+1) \bmod 4})\]



Кільцева топологія (замість лінійного ланцюга) забезпечує рівноправний зв'язок між усіма кубітами, включаючи пару \((q_3, q_0)\), що підвищує ефективність заплутування при тій самій кількості CNOT-гейтів.



Підшар D (варіаційний анзац). Для кожного шару \(l\) використовується окремий набір із 8 навчених параметрів \(\Theta^{(l)} = [\theta_0^{(l)}, \phi_0^{(l)}, \theta_1^{(l)}, \phi_1^{(l)}, \theta_2^{(l)}, \phi_2^{(l)}, \theta_3^{(l)}, \phi_3^{(l)}]\):



\[\forall\,i \in \{0,1,2,3\}:\quad R_y\bigl(\theta_i^{(l)}\bigr)\,|q_i\rangle, \quad R_z\bigl(\phi_i^{(l)}\bigr)\,|q_i\rangle\]



\paragraph{6.7.3. Повна схема}\label{ux43fux43eux432ux43dux430-ux441ux445ux435ux43cux430}



Повна квантова схема складається з \((L + 1)\) повторень шару (п. 6.7.2), за якими слідує вимірювання:



\[U(\mathbf{x}, \Theta) = \prod_{l=0}^{L} \Bigl[\,V(\Theta^{(l)}) \cdot E_{\text{ring}} \cdot ZZ(\mathbf{x}) \cdot W(\mathbf{x})\,\Bigr]\]



де \(W(\mathbf{x})\) --- кодування ознак (підшар A), \(ZZ(\mathbf{x})\) --- ZZ-взаємодії (підшар B), \(E_{\text{ring}}\) --- кільцеве заплутування (підшар C), \(V(\Theta^{(l)})\) --- варіаційний анзац (підшар D).



Вимірювання виконується для всіх 4 кубітів у обчислювальному базисі \(\{|0\rangle,\,|1\rangle\}\).



Витяг ймовірності \(p_{\text{flow}}(t)\). Після виконання схеми \(S\) разів (shots) отримують множину результатів вимірювання у вигляді бітових рядків \(m^{(s)} = (b_0^{(s)},\,b_1^{(s)},\,b_2^{(s)},\,b_3^{(s)}) \in \{0,1\}^4\), де \(b_i^{(s)}\) відповідає результату вимірювання кубіта \(q_i\) у \(s\)-му прогоні. У конкретній реалізації кубіт \(q_0\) визначається як мітковий (readout) кубіт, що кодує клас ``flow'' (\(= 1\)) / ``no-flow'' (\(= 0\)).



Визначається цільова множина результатів для стану потоку:



\[\Omega_{\text{flow}} = \bigl\{\,m \in \{0,1\}^4 :\; b_0 = 1\,\bigr\}\]



Тоді оцінка ймовірності стану потоку обчислюється як маргінальна ймовірність події \(b_0 = 1\):



\[p_{\text{flow}}(t) = \Pr\bigl[b_0 = 1\bigr] \;\approx\; \hat{p}_{\text{flow}}(t) = \frac{1}{S}\sum_{s=1}^{S}\mathbf{1}\bigl[b_0^{(s)} = 1\bigr] = \frac{\displaystyle\sum_{m \in \Omega_{\text{flow}}} \text{count}(m)}{S}\]



Результат \(\hat{p}_{\text{flow}}(t) \in [0,\,1]\).



Параметри \(\Theta\) підбираються (навчання/оптимізація, п. 6.9) на розмічених вікнах з міткою \(\text{FlowLabel}(t) \in \{0,\,1\}\) так, щоб максимізувати правдоподібність події \(b_0 = 1\) для \(\text{FlowLabel} = 1\) та \(b_0 = 0\) для \(\text{FlowLabel} = 0\) (наприклад, за крос-ентропією або іншою функцією втрат), з урахуванням ресурсної вартості \((S,\,D,\,T_{\text{qpu}})\). Таким чином, \(\hat{p}_{\text{flow}}(t)\) є виходом навченого класифікатора, а не довільною частотою.



\paragraph{6.7.4. Параметри за замовчуванням}\label{ux43fux430ux440ux430ux43cux435ux442ux440ux438-ux437ux430-ux437ux430ux43cux43eux432ux447ux443ux432ux430ux43dux43dux44fux43c}



{\def\LTcaptype{none} % do not increment counter

\begin{longtable}[]{@{}

  >{\raggedright\arraybackslash}p{0.5000\linewidth}

  >{\raggedright\arraybackslash}p{0.5000\linewidth}@{}}

\toprule\noalign{}

\begin{minipage}[b]{\linewidth}\raggedright

Параметр

\end{minipage} & \begin{minipage}[b]{\linewidth}\raggedright

Значення

\end{minipage} \\

\midrule\noalign{}

\endhead

\bottomrule\noalign{}

\endlastfoot

Кількість кубітів & \(N_q = 4\) \\

Кількість ознак & \(N_f = 8\) (двоканальне кутове кодування: \(R_y + R_z\)) \\

Кількість шарів повторного завантаження & \(L = 2\) (конфігурований) \\

Загальна кількість навчених параметрів & \(|\Theta| = 24\) \\

Кількість shots & \(S = 1024\) \\

Імітована латентність QPU & 300--2000 мс \\

\end{longtable}

}



\subsubsection{6.8. Обчислення CME у вікні та агрегація}\label{ux43eux431ux447ux438ux441ux43bux435ux43dux43dux44f-cme-ux443-ux432ux456ux43aux43dux456-ux442ux430-ux430ux433ux440ux435ux433ux430ux446ux456ux44f}



CME (обчислювальна ментальна енергія) --- інтегральний операційно визначений показник, що відображає інтенсивність обробки та залученості користувача за EEG-ознаками у заданому часовому вікні \(\Delta\), обчислений на основі спектральної енергії багатоканальних EEG-ознак та модуляції цієї енергії функцією складності задачі і ймовірності стану потоку, оціненої квантовим модулем. Базова одиниця: Вн (\(1\;\text{Вн} \equiv 1\;\text{мкВ}^2 \cdot \text{с}\)).



\paragraph{6.8.1. Миттєва інтенсивність CME}\label{ux43cux438ux442ux442ux454ux432ux430-ux456ux43dux442ux435ux43dux441ux438ux432ux43dux456ux441ux442ux44c-cme}



\[\boxed{\;\text{CME}_{\text{rate}}(t) = \kappa \;\cdot\; E_{\text{band}}(t) \;\cdot\; g\bigl(c(t),\;p_{\text{flow}}(t)\bigr)\;}\]



де \(E_{\text{band}}(t)\) --- зведена спектральна потужність (п. 6.5), одиниці: мкВ²; \(g(c,\,p)\) --- безрозмірна монотонна функція модуляції складності та стану потоку; \(\kappa\) --- безрозмірний коефіцієнт масштабування (п. 6.8.5).



Одиниці \(\text{CME}_{\text{rate}}(t)\): \(\text{мкВ}^2 \equiv \text{Вн/с}\) (оскільки \(1\;\text{Вн} = 1\;\text{мкВ}^2 \cdot \text{с}\), то \(\text{мкВ}^2 = \text{Вн}/\text{с}\)).



\paragraph{6.8.2. Внесок вікна та агрегація по сесії}\label{ux432ux43dux435ux441ux43eux43a-ux432ux456ux43aux43dux430-ux442ux430-ux430ux433ux440ux435ux433ux430ux446ux456ux44f-ux43fux43e-ux441ux435ux441ux456ux457}



Внесок окремого вікна \(t\):



\[\boxed{\;\text{CME}(t) = \text{CME}_{\text{rate}}(t) \;\cdot\; \Delta\;}\]



де \(\Delta\) --- тривалість вікна (с). Одиниці \(\text{CME}(t)\): \(\text{мкВ}^2 \cdot \text{с} \equiv \text{Вн}\).



Сумарна CME за сесію:



\[\text{CME}_{\text{session}} = \sum_{t} \text{CME}(t) = \sum_{t} \text{CME}_{\text{rate}}(t) \cdot \Delta\]



Одиниці: Вн.



\paragraph{6.8.3. Факультативне відображення у систему одиниць SI}\label{ux444ux430ux43aux443ux43bux44cux442ux430ux442ux438ux432ux43dux435-ux432ux456ux434ux43eux431ux440ux430ux436ux435ux43dux43dux44f-ux443-ux441ux438ux441ux442ux435ux43cux443-ux43eux434ux438ux43dux438ux446ux44c-si}



Вн є одиницею \textbf{сигнальної енергії} у сенсі теорії сигналів (\(E = \int |x(t)|^2\,dt\)) і не потребує жодних припущень про опір або імпеданс. Перетворення у Дж є \textbf{необов'язковим} і застосовується лише як апроксимація для порівняння з фізичними одиницями через спрощену електричну модель {[}23{]}:



\[\text{CME}_J(t) = \frac{10^{-12}}{Z_e} \cdot \text{CME}(t)\]



де \(Z_e\) --- імпеданс електрода (Ом), вимірюваний перед сесією (типово 1--10 кОм). Якщо вимір недоступний, використовується нормативне значення \(Z_e = 5\,\text{кОм}\) (IFCN). Одиниці: Дж.



\textbf{Застереження.} Ця формула є наближенням, оскільки: (1) EEG-напруга виникає у кірковій тканині і досягає електрода через об'ємне провідне середовище, яке не є простим резистором; (2) \(Z_e\) є частотно-залежним і скалярне значення є усередненням. Основним результатом системи є CME у Вн; відображення у Дж слугує виключно довідковим орієнтиром.



Для зручності інтерпретації та візуалізації визначається похідна безрозмірна шкала:



\[\text{CME}_{\text{index}}(t) = \frac{\text{CME}_{\text{rate}}(t)}{\text{CME}_{\text{rate,max}}^{(\text{cal})}} \times 100\]



де \(\text{CME}_{\text{rate,max}}^{(\text{cal})}\) --- максимальне значення \(\text{CME}_{\text{rate}}\), визначене під час калібрування. \(\text{CME}_{\text{index}}(t) \in [0,\,100]\) є безрозмірною шкалою відображення і не є базовою одиницею CME.



\paragraph{6.8.4. Функція модуляції}\label{ux444ux443ux43dux43aux446ux456ux44f-ux43cux43eux434ux443ux43bux44fux446ux456ux457}



У загальному вигляді \(g(c,\,p)\) є конфігурованою безрозмірною функцією від нормованої складності задачі \(c \in [0,\,1]\) та ймовірності стану потоку \(p \in [0,\,1]\). У переважному варіанті реалізації використовується лінійно-білінійна форма:



\[g(c,\,p) = \lambda_1 \cdot c + \lambda_2 \cdot p + \lambda_3 \cdot c \cdot p\]



де \(\lambda_1, \lambda_2, \lambda_3 \geq 0\) --- безрозмірні вагові коефіцієнти. Значення \(g(c,\,p) \in [0,\,\lambda_1 + \lambda_2 + \lambda_3]\); за потреби \(g(c,\,p)\) може бути обмежена до інтервалу \([0,\,g_{\max}]\) для стабілізації масштабу. Допускаються й інші форми функції модуляції (наприклад, степенева \(g(c,p) = c^\alpha \cdot p^\beta\)), за умови збереження безрозмірності та монотонності по кожному аргументу.



Значення коефіцієнтів за замовчуванням:



{\def\LTcaptype{none} % do not increment counter

\begin{longtable}[]{@{}lll@{}}

\toprule\noalign{}

Коефіцієнт & Значення & Призначення \\

\midrule\noalign{}

\endhead

\bottomrule\noalign{}

\endlastfoot

\(\lambda_1\) & 0.5 & складність задачі \\

\(\lambda_2\) & 0.5 & ймовірність стану потоку \\

\(\lambda_3\) & 0.5 & добуток складності та ймовірності стану потоку \\

\end{longtable}

}



\paragraph{\texorpdfstring{6.8.5. Коефіцієнт масштабування \(\kappa\)}{6.8.5. Коефіцієнт масштабування \textbackslash kappa}}\label{ux43aux43eux435ux444ux456ux446ux456ux454ux43dux442-ux43cux430ux441ux448ux442ux430ux431ux443ux432ux430ux43dux43dux44f-kappa}



\(\kappa\) є безрозмірним коефіцієнтом, що забезпечує узгодження масштабу CME між різними конфігураціями, сесіями та користувачами. Значення \(\kappa\) визначається одним із способів (у порядку пріоритету):



\begin{enumerate}

\def\labelenumi{\arabic{enumi}.}

\tightlist

\item

  За персональним профілем: \(\kappa_{\text{user}}\), обчислений за попередніми сесіями користувача та збережений у сховищі. Забезпечує міжсесійну порівнянність для одного користувача.

\item

  За калібрувальним періодом: перші \(N_{\text{cal}}\) вікон поточної сесії (наприклад, перші 2--5 хвилин baseline). Коефіцієнт обчислюється як:

\end{enumerate}



\[\kappa = \frac{C_{\text{target}}}{\displaystyle\max_{t \in T_{\text{cal}}}\;\text{CME}_{\text{rate,raw}}(t)}\]



де \(C_{\text{target}}\) --- конфігурований цільовий максимум шкали (наприклад, \(C_{\text{target}} = 100\) для шкали \([0,\,100]\)), \(\text{CME}_{\text{rate,raw}}(t) = E_{\text{band}}(t) \cdot g(c(t),\,p_{\text{flow}}(t))\) --- некаліброване значення інтенсивності. \(\kappa\) фіксується після калібрування та не змінюється протягом сесії.



\begin{enumerate}

\def\labelenumi{\arabic{enumi}.}

\setcounter{enumi}{2}

\tightlist

\item

  За фіксованою константою: якщо персональний профіль та калібрувальні дані відсутні, використовується \(\kappa_{\text{default}}\), визначена емпірично для типових діапазонів спектральних потужностей.

\end{enumerate}



\paragraph{6.8.6. Метрики сесії}\label{ux43cux435ux442ux440ux438ux43aux438-ux441ux435ux441ux456ux457}



Додатково обчислюються такі агреговані показники:



{\def\LTcaptype{none} % do not increment counter

\begin{longtable}[]{@{}

  >{\raggedright\arraybackslash}p{0.3750\linewidth}

  >{\raggedright\arraybackslash}p{0.2500\linewidth}

  >{\raggedright\arraybackslash}p{0.3750\linewidth}@{}}

\toprule\noalign{}

\begin{minipage}[b]{\linewidth}\raggedright

Метрика

\end{minipage} & \begin{minipage}[b]{\linewidth}\raggedright

Опис

\end{minipage} & \begin{minipage}[b]{\linewidth}\raggedright

Одиниці

\end{minipage} \\

\midrule\noalign{}

\endhead

\bottomrule\noalign{}

\endlastfoot

\(N_{\text{total}}\), \(N_{\text{flow}}\) & загальна кількість вікон та кількість вікон у стані потоку & шт. \\

\(\text{FlowShare}\) & \(N_{\text{flow}} / N_{\text{total}}\) & безрозм. \\

\(\overline{\text{CME}}\) & середнє CME по вікнах & Вн \\

\(\text{CME}_{\max}\) & максимальне CME & Вн \\

\(\text{LongestFlowStreak}\) & тривалість найдовшої безперервної серії стану потоку & вікон \\

Періоди стану потоку & часові мітки початку/кінця та середні \(\text{CME}\), \(p_{\text{flow}}\) & --- \\

\end{longtable}

}



\paragraph{6.8.7. Детектування стану потоку}\label{ux434ux435ux442ux435ux43aux442ux443ux432ux430ux43dux43dux44f-ux441ux442ux430ux43dux443-ux43fux43eux442ux43eux43aux443}



Вікно вважається вікном у стані потоку, якщо:



\[p_{\text{flow}}(t) \;\geq\; \text{FlowThreshold}\]



де \(\text{FlowThreshold}\) --- конфігурований пороговий коефіцієнт (за замовчуванням 0.85). Поріг може визначатися:



\begin{enumerate}

\def\labelenumi{\arabic{enumi}.}

\tightlist

\item

  як фіксоване значення з конфігурації;

\item

  на основі калібрування (наприклад, як percentile \(p_{90}\) розподілу \(p_{\text{flow}}\) під час калібрувального baseline-періоду);

\item

  індивідуально для кожного користувача на основі персонального профілю.

\end{enumerate}



Оскільки \(\hat{p}_{\text{flow}}(t)\) оцінюється за скінченним \(S\), дисперсія оцінки зменшується зі зростанням \(S\); за потреби рішення ``flow'' може прийматись за умови \(\hat{p}_{\text{flow}}(t) \geq \text{FlowThreshold}\) або за нижньою межею довірчого інтервалу.



\subsubsection{6.9. Метаевристична оптимізація квантових параметрів та ресурсів}\label{ux43cux435ux442ux430ux435ux432ux440ux438ux441ux442ux438ux447ux43dux430-ux43eux43fux442ux438ux43cux456ux437ux430ux446ux456ux44f-ux43aux432ux430ux43dux442ux43eux432ux438ux445-ux43fux430ux440ux430ux43cux435ux442ux440ux456ux432-ux442ux430-ux440ux435ux441ux443ux440ux441ux456ux432}



Для підвищення робастності та/або якості оцінювання застосовується метаевристичний оптимізатор (190), реалізований як фоновий програмний сервіс. Підхід спирається на результати мета-оптимізації квантових ресурсів {[}12{]} та враховує обмеження shots і латентності, відомі із задач планування квантових завдань {[}13, 20--22{]}.



Вектор оптимізації визначається як \(\mathbf{u} = (\Theta,\;S,\;D)\), де \(\Theta \in \mathbb{R}^{8(L+1)}\) --- параметри варіаційних шарів квантової схеми (24 числа для \(L = 2\)), \(S\) --- кількість shots, \(D\) --- глибина схеми.



\paragraph{6.9.1. Цільова функція}\label{ux446ux456ux43bux44cux43eux432ux430-ux444ux443ux43dux43aux446ux456ux44f}



\[\boxed{\;J(\mathbf{u}) = \mathcal{L}_{\text{val}}(\mathbf{u}) + \xi_S \cdot S + \xi_D \cdot D + \xi_T \cdot T_{\text{qpu}}(\mathbf{u})\;}\]



де \(\mathcal{L}_{\text{val}}(\mathbf{u})\) --- помилка/нестійкість оцінювання на калібрувальних/валідаційних даних; \(T_{\text{qpu}}(\mathbf{u})\) --- виміряний або оцінений час виконання на QPU/симуляторі; \(\xi_S,\;\xi_D,\;\xi_T \geq 0\) --- вагові коефіцієнти штрафу за ресурси (обрано відмінно від \(\alpha,\beta,\gamma\) --- назв частотних діапазонів EEG).



\paragraph{6.9.2. Алгоритм навчання}\label{ux430ux43bux433ux43eux440ux438ux442ux43c-ux43dux430ux432ux447ux430ux43dux43dux44f}



\begin{enumerate}

\def\labelenumi{\arabic{enumi}.}

\item

  Ініціалізація популяції: генерується \(N\) кандидатів (за замовчуванням \(N = 5\)), кожен із випадковими параметрами \(\Theta \in [0,\,2\pi]^{24}\) (по 8 параметрів на кожен із \(L + 1 = 3\) шарів).

\item

  Цикл поколінь (за замовчуванням 10 поколінь):



  \begin{enumerate}

  \def\labelenumii{\alph{enumii}.}

  \tightlist

  \item

    Для кожного кандидата з валідаційного набору (зібрані та розмічені вікна EEG з набору даних, п. 6.10) обирається пакет (batch) вікон; для кожного вікна виконується квантовий інференс через QPU/симулятор із параметрами кандидата; обчислюється фітнес-значення як метрика якості на валідаційному наборі, наприклад бінарна крос-ентропія або F1-score: \(\text{fitness} = \mathcal{L}_{\text{val}}(\Theta_k)\), де \(\mathcal{L}_{\text{val}}\) --- усереднена по пакету метрика, що порівнює \(p_{\text{flow}}(t)\) з мітками \(\text{FlowLabel}(t)\). При стохастичній оцінці (обумовленій квантовим шумом при кінцевому \(S\)) виконується усереднення по \(R \geq 3\) повторах.

  \item

    Після оцінки всіх кандидатів виконується еволюція популяції.

  \end{enumerate}

\item

  Еволюція популяції складається з трьох операцій: (а) відбір --- 50 \% найкращих кандидатів виживають; (б) кросовер --- нащадки формуються вибором параметрів від двох батьків (uniform crossover); (в) мутація --- кожен параметр мутує з ймовірністю \(p_{\text{mut}}\) та величиною \(\pm 0.25\) (рівномірний розподіл).

\end{enumerate}



{\def\LTcaptype{none} % do not increment counter

\begin{longtable}[]{@{}lc@{}}

\toprule\noalign{}

Алгоритм & \(p_{\text{mut}}\) \\

\midrule\noalign{}

\endhead

\bottomrule\noalign{}

\endlastfoot

PSO & 0.05 \\

GA & 0.10 \\

ACO & 0.15 \\

SA & 0.20 \\

\end{longtable}

}



\begin{enumerate}

\def\labelenumi{\arabic{enumi}.}

\setcounter{enumi}{3}

\tightlist

\item

  Збереження найкращої моделі: після завершення навчання параметри найкращого кандидата зберігаються у сховище і позначаються як активна модель. Попередні активні моделі деактивуються. Під час інференсу система автоматично завантажує навчені параметри для використання у варіаційній частині квантової схеми.

\end{enumerate}



\paragraph{6.9.3. Підтримувані метаевристики}\label{ux43fux456ux434ux442ux440ux438ux43cux443ux432ux430ux43dux456-ux43cux435ux442ux430ux435ux432ux440ux438ux441ux442ux438ux43aux438}



{\def\LTcaptype{none} % do not increment counter

\begin{longtable}[]{@{}

  >{\raggedright\arraybackslash}p{0.4800\linewidth}

  >{\raggedright\arraybackslash}p{0.2800\linewidth}

  >{\raggedright\arraybackslash}p{0.2400\linewidth}@{}}

\toprule\noalign{}

\begin{minipage}[b]{\linewidth}\raggedright

Позначення

\end{minipage} & \begin{minipage}[b]{\linewidth}\raggedright

Назва

\end{minipage} & \begin{minipage}[b]{\linewidth}\raggedright

Опис

\end{minipage} \\

\midrule\noalign{}

\endhead

\bottomrule\noalign{}

\endlastfoot

GA & Genetic Algorithm & генетичний алгоритм із відбором, кросовером та мутацією \\

PSO & Particle Swarm Optimization & оптимізація роєм частинок \\

ACO & Ant Colony Optimization & мурашиний алгоритм \\

SA & Simulated Annealing & імітація відпалу \\

\end{longtable}

}



\subsubsection{6.10. Модуль класичного нейромережевого аналізу та набір даних}\label{ux43cux43eux434ux443ux43bux44c-ux43aux43bux430ux441ux438ux447ux43dux43eux433ux43e-ux43dux435ux439ux440ux43eux43cux435ux440ux435ux436ux435ux432ux43eux433ux43e-ux430ux43dux430ux43bux456ux437ux443-ux442ux430-ux43dux430ux431ux456ux440-ux434ux430ux43dux438ux445}



\paragraph{6.10.1. Збір та збереження набору даних EEG}\label{ux437ux431ux456ux440-ux442ux430-ux437ux431ux435ux440ux435ux436ux435ux43dux43dux44f-ux43dux430ux431ux43eux440ux443-ux434ux430ux43dux438ux445-eeg}



Паралельно з обчисленням CME система забезпечує автоматичний збір і збереження набору даних EEG-ознак. Для кожного вікна \(t\) створюється запис, що містить:



{\def\LTcaptype{none} % do not increment counter

\begin{longtable}[]{@{}

  >{\raggedright\arraybackslash}p{0.7391\linewidth}

  >{\raggedright\arraybackslash}p{0.2609\linewidth}@{}}

\toprule\noalign{}

\begin{minipage}[b]{\linewidth}\raggedright

Складова запису

\end{minipage} & \begin{minipage}[b]{\linewidth}\raggedright

Опис

\end{minipage} \\

\midrule\noalign{}

\endhead

\bottomrule\noalign{}

\endlastfoot

Спектральні потужності & повний 20-канальний вектор \(P_{\text{band}}(ch, t)\) по 4 електродах у 5 діапазонах \\

Контекстні індекси & складність задачі \(c(t)\) та індикатор якості сигналу \(q(t)\) \\

Мітки стану потоку (опційно) & бінарна мітка та ймовірність, що присвоюються після маркування \\

Маркер дії & часовий інтервал із зазначенням типу дії користувача \\

\end{longtable}

}



Модуль асинхронного збереження (205) отримує дані через обмежену чергу і записує їх у пакетному режимі, що гарантує неблокуючий характер запису для основного потоку обробки.



При записі кожного вікна автоматично виконується:



\begin{enumerate}

\def\labelenumi{\arabic{enumi}.}

\item

  Пошук маркера дії --- знаходження маркера дії, часовий інтервал якого покриває мітку часу вікна, з пріоритетом найдовшого інтервалу.

\item

  Евристичне маркування --- обчислення початкової мітки стану потоку за евристикою фронтального \(\alpha/\theta\)-відношення:

\end{enumerate}



\[r(t) = \frac{P_\alpha(\text{AF7}, t) + P_\alpha(\text{AF8}, t)}{P_\theta(\text{AF7}, t) + P_\theta(\text{AF8}, t)}\]



\[\text{FlowLabel}_{\text{heur}}(t) = \begin{cases}1, & r(t) \geq 1.0 \\ 0, & r(t) < 1.0\end{cases}, \qquad p_{\text{flow}}^{\text{heur}}(t) = \text{clamp}\!\left(\frac{r(t)}{2},\;0,\;1\right)\]



Ця евристика базується на дослідженнях, що показують зв'язок між підвищеним фронтальним \(\alpha/\theta\)-відношенням та станом зосередженості {[}4, 6{]}. Евристична мітка використовується як початкове слабке маркування (weak labeling) і може бути замінена або уточнена вручну, результатами зовнішнього класифікатора чи експериментального протоколу.



\textbf{Валідація міток стану потоку.} Для верифікації автоматичних міток передбачається паралельне застосування визнаних психометричних шкал: Flow State Scale (FSS-2) {[}24{]} та/або Flow-Kurzskala (FKS) {[}25{]}. Процедура полягає у заповненні відповідного опитувальника після завершення сесії; отриманий сумарний бал порівнюється із середнім значенням \(p_{\text{flow}}(t)\) та часткою вікон у стані потоку за сесію. Кореляція між суб'єктивним балом та об'єктивним показником CME використовується для калібрування порогу FlowThreshold та оцінки якості маркування.



\paragraph{6.10.2. Класичний нейромережевий аналіз}\label{ux43aux43bux430ux441ux438ux447ux43dux438ux439-ux43dux435ux439ux440ux43eux43cux435ux440ux435ux436ux435ux432ux438ux439-ux430ux43dux430ux43bux456ux437}



Модуль класичного нейромережевого аналізу (195) реалізований як окремий програмний сервіс з навченою нейронною мережею, що приймає 22-ознаковий вхідний вектор:



\[\mathbf{f}_t = \Bigl[\,P_\delta^{\text{TP9}},\;P_\theta^{\text{TP9}},\;\ldots,\;P_\gamma^{\text{TP10}},\;c(t),\;q(t)\,\Bigr] \in \mathbb{R}^{22}\]



та повертає: (а) ймовірність стану потоку \(p_{\text{flow}}^{\text{NN}}(t) \in [0,\,1]\); (б) бінарну мітку \(\text{FlowLabel}_{\text{NN}}(t) = \mathbb{1}\bigl[p_{\text{flow}}^{\text{NN}}(t) \geq 0.5\bigr]\).



Серверний програмний модуль ітеративно передає немарковані вікна до класифікатора, отримує мітки та зберігає їх у сховищі. Це дозволяє:



\begin{enumerate}

\def\labelenumi{\arabic{enumi}.}

\tightlist

\item

  масово маркувати зібрані дані для подальшого навчання квантової моделі;

\item

  порівнювати якість класичної та квантової моделей на тому самому наборі;

\item

  формувати гібридний інференс, де квантова та класична моделі доповнюють одна одну.

\end{enumerate}



\paragraph{6.10.3. Маркери дій}\label{ux43cux430ux440ux43aux435ux440ux438-ux434ux456ux439}



Для побудови розміченого набору даних система підтримує маркери дій --- часові інтервали \([t_{\text{start}},\,t_{\text{end}}]\) із типом дії (``coding'', ``reading'', ``rest'', ``meditation'' тощо). Маркери створюються через програмний інтерфейс або автоматично за допомогою зовнішніх інтеграцій. Кожне вікно EEG прив'язується до маркера дії, часовий інтервал якого покриває мітку часу вікна, що дозволяє аналізувати нейрофізіологічний стан у контексті конкретної діяльності.



\paragraph{6.10.4. Режими інференсу}\label{ux440ux435ux436ux438ux43cux438-ux456ux43dux444ux435ux440ux435ux43dux441ux443}



Система підтримує три режими інференсу для оцінювання стану потоку у кожному вікні:



{\def\LTcaptype{none} % do not increment counter

\begin{longtable}[]{@{}

  >{\raggedright\arraybackslash}p{0.1667\linewidth}

  >{\raggedright\arraybackslash}p{0.1905\linewidth}

  >{\raggedright\arraybackslash}p{0.3810\linewidth}

  >{\raggedright\arraybackslash}p{0.2619\linewidth}@{}}

\toprule\noalign{}

\begin{minipage}[b]{\linewidth}\raggedright

Режим

\end{minipage} & \begin{minipage}[b]{\linewidth}\raggedright

Модель

\end{minipage} & \begin{minipage}[b]{\linewidth}\raggedright

Вхідний вектор

\end{minipage} & \begin{minipage}[b]{\linewidth}\raggedright

Результат

\end{minipage} \\

\midrule\noalign{}

\endhead

\bottomrule\noalign{}

\endlastfoot

Квантовий & VQC (4 кубіти) & \(\mathbf{x}_t^{(q)} \in \mathbb{R}^8\) & \(p_{\text{flow}}(t)\) --- CME \\

Класичний & NN (класифікатор стану потоку) & \(\mathbf{f}_t \in \mathbb{R}^{22}\) & \(p_{\text{flow}}^{\text{NN}}(t)\) \\

Гібридний & VQC + NN & \(\mathbf{x}_t^{(q)}\), \(\mathbf{f}_t\) & зважена комбінація \\

\end{longtable}

}



У гібридному режимі підсумкова ймовірність стану потоку може обчислюватись як:



\[p_{\text{flow}}^{\text{hybrid}}(t) = \mu \cdot p_{\text{flow}}(t) + (1 - \mu) \cdot p_{\text{flow}}^{\text{NN}}(t)\]



де \(\mu \in [0,\,1]\) --- конфігурований коефіцієнт довіри до квантової моделі.



\subsubsection{6.11. Варіанти розгортання}\label{ux432ux430ux440ux456ux430ux43dux442ux438-ux440ux43eux437ux433ux43eux440ux442ux430ux43dux43dux44f}



Система підтримує щонайменше три варіанти архітектури:



\begin{enumerate}

\def\labelenumi{\Alph{enumi})}

\item

  Монолітна архітектура. Усі модулі (попередня обробка, звернення до QPU, обчислення CME, запис у сховище) виконуються синхронно в одному процесі/сервісі. Послідовність операцій: валідація --- попередня обробка --- виклик QPU --- обчислення CME --- запис у сховище --- відповідь клієнту.

\item

  Синхронна мікросервісна композиція. Програмний інтерфейс викликає окремий сервіс попередньої обробки через мережевий протокол, потім квантовий модуль, потім записує результат у сховище. Між сервісами може виникати мережева затримка.

\item

  Брокеризована архітектура. Програмний інтерфейс ставить запити у чергу повідомлень (брокер). Робочі вузли зчитують запити з черги та виконують повний цикл обробки (попередня обробка --- QPU --- CME --- сховище). Результат повертається клієнту після завершення обробки.

\end{enumerate}



Зазначені варіанти не змінюють математичної суті винаходу, а описують різні режими реалізації єдиного задуму, що дозволяє масштабувати систему від одного вузла до розподіленого кластера.



\subsubsection{6.12. Контейнеризація та розгортання}\label{ux43aux43eux43dux442ux435ux439ux43dux435ux440ux438ux437ux430ux446ux456ux44f-ux442ux430-ux440ux43eux437ux433ux43eux440ux442ux430ux43dux43dux44f}



Система може розгортатися у контейнерному середовищі (наприклад, Docker Compose) із такими компонентами:



{\def\LTcaptype{none} % do not increment counter

\begin{longtable}[]{@{}

  >{\raggedright\arraybackslash}p{0.4583\linewidth}

  >{\raggedright\arraybackslash}p{0.5417\linewidth}@{}}

\toprule\noalign{}

\begin{minipage}[b]{\linewidth}\raggedright

Компонент

\end{minipage} & \begin{minipage}[b]{\linewidth}\raggedright

Призначення

\end{minipage} \\

\midrule\noalign{}

\endhead

\bottomrule\noalign{}

\endlastfoot

Реляційна база даних & зберігання сесій, результатів, моделей (SQL Server або PostgreSQL) \\

Сервіс квантового модуля & виконання квантового інференсу \\

Сервіс класичного класифікатора & нейромережевий аналіз стану потоку \\

Серверний програмний інтерфейс & Web API для координації та доступу \\

Веб-інтерфейс & візуалізація та керування \\

\end{longtable}

}



Усі компоненти з'єднані через мережу та взаємодіють через HTTP. Пристрій стримінгу (115) може працювати окремо на одноплатному комп'ютері та підключатися до серверного модуля через мережу.



\subsubsection{6.13. Взаємодія компонентів системи}\label{ux432ux437ux430ux454ux43cux43eux434ux456ux44f-ux43aux43eux43cux43fux43eux43dux435ux43dux442ux456ux432-ux441ux438ux441ux442ux435ux43cux438}



Ключові аспекти взаємодії компонентів:



{\def\LTcaptype{none} % do not increment counter

\begin{longtable}[]{@{}

  >{\raggedright\arraybackslash}p{0.1071\linewidth}

  >{\raggedright\arraybackslash}p{0.3214\linewidth}

  >{\raggedright\arraybackslash}p{0.3571\linewidth}

  >{\raggedright\arraybackslash}p{0.2143\linewidth}@{}}

\toprule\noalign{}

\begin{minipage}[b]{\linewidth}\raggedright

№

\end{minipage} & \begin{minipage}[b]{\linewidth}\raggedright

Зв'язок

\end{minipage} & \begin{minipage}[b]{\linewidth}\raggedright

Протокол

\end{minipage} & \begin{minipage}[b]{\linewidth}\raggedright

Опис

\end{minipage} \\

\midrule\noalign{}

\endhead

\bottomrule\noalign{}

\endlastfoot

1 & Пристрій реєстрації (110) --- пристрій стримінгу (115) & Bluetooth & передача сирих даних EEG \\

2 & Пристрій стримінгу (115) --- сервер (120) & OSC / WebSocket / SignalR / HTTP & буферизація, попередня сегментація та передача вікон \\

3 & Сервер --- квантовий бекенд (170) & HTTP & передача ознак для інференсу, отримання \(p_{\text{flow}}(t)\) \\

4 & Сервер --- модуль збереження (205) & паралельна черга & асинхронний запис сирих даних незалежно від результатів CME \\

5 & Сервер --- класичний класифікатор (195) & HTTP & масове маркування вікон, збереження міток стану потоку \\

6 & Метаевристика (190) --- QPU & внутрішній & фонова оптимізація параметрів \(\Theta\) квантової схеми \\

7 & Веб-інтерфейс (210) --- сервер & WebSocket / REST API & стримінг CME у реальному часі, керування сесіями, аналіз, експорт \\

\end{longtable}

}



\begin{center}\rule{0.5\linewidth}{0.5pt}\end{center}



\subsection{7. Опис, що підтверджує можливість реалізації (приклад виконання)}\label{ux43eux43fux438ux441-ux449ux43e-ux43fux456ux434ux442ux432ux435ux440ux434ux436ux443ux454-ux43cux43eux436ux43bux438ux432ux456ux441ux442ux44c-ux440ux435ux430ux43bux456ux437ux430ux446ux456ux457-ux43fux440ux438ux43aux43bux430ux434-ux432ux438ux43aux43eux43dux430ux43dux43dux44f}



\subsubsection{7.1. Приклад вхідних даних}\label{ux43fux440ux438ux43aux43bux430ux434-ux432ux445ux456ux434ux43dux438ux445-ux434ux430ux43dux438ux445}



Приклад вхідного запису для одного вікна \(\Delta\) у форматі CSV від пристрою стримінгу або експорту з системи:



{\def\LTcaptype{none} % do not increment counter

\begin{longtable}[]{@{}

  >{\raggedright\arraybackslash}p{0.1463\linewidth}

  >{\centering\arraybackslash}p{0.1220\linewidth}

  >{\centering\arraybackslash}p{0.1220\linewidth}

  >{\centering\arraybackslash}p{0.1220\linewidth}

  >{\centering\arraybackslash}p{0.1220\linewidth}

  >{\centering\arraybackslash}p{0.1220\linewidth}

  >{\centering\arraybackslash}p{0.1220\linewidth}

  >{\centering\arraybackslash}p{0.1220\linewidth}@{}}

\toprule\noalign{}

\begin{minipage}[b]{\linewidth}\raggedright

Time

\end{minipage} & \begin{minipage}[b]{\linewidth}\centering

\(P_\delta^{\text{AF7}}\)

\end{minipage} & \begin{minipage}[b]{\linewidth}\centering

\(P_\theta^{\text{AF7}}\)

\end{minipage} & \begin{minipage}[b]{\linewidth}\centering

\(P_\alpha^{\text{AF7}}\)

\end{minipage} & \begin{minipage}[b]{\linewidth}\centering

\(P_\beta^{\text{AF7}}\)

\end{minipage} & \begin{minipage}[b]{\linewidth}\centering

\(P_\gamma^{\text{AF7}}\)

\end{minipage} & \begin{minipage}[b]{\linewidth}\centering

\(Q\)

\end{minipage} & \begin{minipage}[b]{\linewidth}\centering

\(c\)

\end{minipage} \\

\midrule\noalign{}

\endhead

\bottomrule\noalign{}

\endlastfoot

2025-09-03 19:51:54.599 & 1.0795 & 0.3091 & 0.3864 & 0.1206 & 0.0578 & 0.93 & 0.62 \\

\end{longtable}

}



Позначення: Time --- часова мітка; \(P_\delta^{\text{AF7}}\), \(P_\theta^{\text{AF7}}\), \(P_\alpha^{\text{AF7}}\), \(P_\beta^{\text{AF7}}\), \(P_\gamma^{\text{AF7}}\) --- спектральні потужності у діапазонах \(\delta/\theta/\alpha/\beta/\gamma\) для електроду AF7; \(Q\) --- індикатор якості сигналу; \(c\) --- індекс складності.



\subsubsection{7.2. Приклад обчислення (крок за кроком)}\label{ux43fux440ux438ux43aux43bux430ux434-ux43eux431ux447ux438ux441ux43bux435ux43dux43dux44f-ux43aux440ux43eux43a-ux437ux430-ux43aux440ux43eux43aux43eux43c}



Крок 1 (сегментація). Модуль сегментації формує вікно \(\Delta = 5\) с із часовими мітками початку та кінця.



Крок 2 (виділення ознак). Спектральні потужності по всіх 4 електродах (Фіг. 4):



{\def\LTcaptype{none} % do not increment counter

\begin{longtable}[]{@{}lllll@{}}

\toprule\noalign{}

Електрод & \(P_\delta\) & \(P_\theta\) & \(P_\alpha\) & \(P_\beta\) \\

\midrule\noalign{}

\endhead

\bottomrule\noalign{}

\endlastfoot

TP9 & 0.9821 & 0.2756 & 0.4012 & 0.1345 \\

AF7 & 1.0795 & 0.3091 & 0.3864 & 0.1206 \\

AF8 & 1.1234 & 0.2987 & 0.3956 & 0.1289 \\

TP10 & 0.9567 & 0.2834 & 0.4123 & 0.1398 \\

\end{longtable}

}



Крок 3 (обчислення \(E_{\text{band}}(t)\)). Ваги за замовчуванням (п. 6.5): \(w_\delta = 0.5\), \(w_\theta = 1.0\), \(w_\alpha = 1.0\), \(w_\beta = 0.3\), \(w_\gamma = 0\). Обчислення по кожному електроду та сумування:



\[\begin{aligned}

E_{\text{TP9}} &= 0.5 \times 0.9821 + 1.0 \times 0.2756 + 1.0 \times 0.4012 + 0.3 \times 0.1345 = 1.2082 \\

E_{\text{AF7}} &= 0.5 \times 1.0795 + 1.0 \times 0.3091 + 1.0 \times 0.3864 + 0.3 \times 0.1206 = 1.2714 \\

E_{\text{AF8}} &= 0.5 \times 1.1234 + 1.0 \times 0.2987 + 1.0 \times 0.3956 + 0.3 \times 0.1289 = 1.2947 \\

E_{\text{TP10}} &= 0.5 \times 0.9567 + 1.0 \times 0.2834 + 1.0 \times 0.4123 + 0.3 \times 0.1398 = 1.2160 \\[4pt]

E_{\text{band}}(t) &= 1.2082 + 1.2714 + 1.2947 + 1.2160 = 4.9903 \;\text{мкВ}^2

\end{aligned}\]



Крок 4 (формування \(\mathbf{x}_t^{(q)}\) та квантовий інференс). Зведений 8-ознаковий вектор \(\mathbf{x}_t^{(q)} = [\,0.52,\;0.15,\;0.19,\;0.06,\;0.03,\;-0.12,\;0.49,\;0.62\,]\) (нормалізовані середні потужності \(\overline{P}_{\delta/\theta/\alpha/\beta/\gamma}\), фронтальна \(\alpha\)-асиметрія, індекс залученості \(\beta/\theta\), складність \(c(t)\)) подається у квантову схему. Після виконання \(S = 1024\) shots квантова схема повертає:



\[p_{\text{flow}}(t) = 0.623\]



При роботі на симуляторі для моделювання шуму реального QPU опційно може додаватися збурення \(\varepsilon \sim U(-0.02,\,0.02)\) до \(p_{\text{flow}}\). У продуктивному режимі (реальний QPU) збурення не застосовується.



Крок 5 (обчислення \(\text{CME}(t)\)). З індексом складності \(c(t) = 0.62\):



\[\begin{aligned}

g(c,\,p) &= 0.5 \times 0.62 + 0.5 \times 0.623 + 0.5 \times 0.62 \times 0.623 \\

&= 0.31 + 0.3115 + 0.19313 \\

&= 0.81463

\end{aligned}\]



\[\text{CME}_{\text{rate,raw}}(t) = 4.9903 \times 0.81463 = 4.0652 \;\text{Вн/с}\]



У цьому прикладі приймається \(\kappa = 1\) (некалібрований режим, п. 6.8.5):



\[\text{CME}_{\text{rate}}(t) = 1 \times 4.0652 = 4.0652 \;\text{Вн/с}\]



\[\text{CME}(t) = \text{CME}_{\text{rate}}(t) \times \Delta = 4.0652 \;\text{Вн/с} \times 5.0 \;\text{с} = 20.326 \;\text{Вн}\]



Для безрозмірної шкали відображення (за калібрувальним максимумом \(\text{CME}_{\text{rate,max}}^{(\text{cal})}\)): \(\text{CME}_{\text{index}}(t) = 4.0652 / \text{CME}_{\text{rate,max}}^{(\text{cal})} \times 100\).



Крок 6 (агрегація).



\[\text{CME}_{\text{session}} = \sum_{t} \text{CME}(t) \;\text{(Вн)} \quad\text{для всіх вікон сесії.}\]



\subsubsection{7.3. Приклад формату API-запиту}\label{ux43fux440ux438ux43aux43bux430ux434-ux444ux43eux440ux43cux430ux442ux443-api-ux437ux430ux43fux438ux442ux443}



Обмін даними між модулями здійснюється через структуровані повідомлення. Нижче наведено ілюстративний приклад у форматі JSON/REST; конкретний формат та протокол передачі не є обов'язковими ознаками винаходу.



Запит до ендпоінту онлайн-інференсу:



\begin{Shaded}

\begin{Highlighting}[]

\ErrorTok{POST} \ErrorTok{/api/inference/cme}

\ErrorTok{Content{-}Type:} \ErrorTok{application/json}



\FunctionTok{\{}

    \DataTypeTok{"sessionId"}\FunctionTok{:} \StringTok{"a1b2c3d4{-}e5f6{-}7890{-}abcd{-}ef1234567890"}\FunctionTok{,}

    \DataTypeTok{"windowId"}\FunctionTok{:} \StringTok{"w\_001"}\FunctionTok{,}

    \DataTypeTok{"features"}\FunctionTok{:} \OtherTok{[}\FloatTok{0.52}\OtherTok{,} \FloatTok{0.15}\OtherTok{,} \FloatTok{0.19}\OtherTok{,} \FloatTok{0.06}\OtherTok{,} \FloatTok{0.03}\OtherTok{,} \FloatTok{{-}0.12}\OtherTok{,} \FloatTok{0.49}\OtherTok{,} \FloatTok{0.62}\OtherTok{]}\FunctionTok{,}

    \DataTypeTok{"taskDifficulty"}\FunctionTok{:} \FloatTok{0.62}

\FunctionTok{\}}

\end{Highlighting}

\end{Shaded}



Відповідь:



\begin{Shaded}

\begin{Highlighting}[]

\FunctionTok{\{}

    \DataTypeTok{"cmeVn"}\FunctionTok{:} \FloatTok{20.326}\FunctionTok{,}

    \DataTypeTok{"cmeIndex"}\FunctionTok{:} \FloatTok{73.45}\FunctionTok{,}

    \DataTypeTok{"pFlow"}\FunctionTok{:} \FloatTok{0.623}\FunctionTok{,}

    \DataTypeTok{"shotsUsed"}\FunctionTok{:} \DecValTok{1024}\FunctionTok{,}

    \DataTypeTok{"depth"}\FunctionTok{:} \DecValTok{8}\FunctionTok{,}

    \DataTypeTok{"qpuLatencyMs"}\FunctionTok{:} \DecValTok{1456}\FunctionTok{,}

    \DataTypeTok{"totalLatencyMs"}\FunctionTok{:} \DecValTok{1589}

\FunctionTok{\}}

\end{Highlighting}

\end{Shaded}



\subsubsection{7.4. Приклад запуску навчання}\label{ux43fux440ux438ux43aux43bux430ux434-ux437ux430ux43fux443ux441ux43aux443-ux43dux430ux432ux447ux430ux43dux43dux44f}



Запит на створення завдання метаевристичної оптимізації:



\begin{Shaded}

\begin{Highlighting}[]

\ErrorTok{POST} \ErrorTok{/api/training/start}

\ErrorTok{Content{-}Type:} \ErrorTok{application/json}



\FunctionTok{\{}

    \DataTypeTok{"algorithm"}\FunctionTok{:} \StringTok{"genetic"}\FunctionTok{,}

    \DataTypeTok{"totalGenerations"}\FunctionTok{:} \DecValTok{10}

\FunctionTok{\}}

\end{Highlighting}

\end{Shaded}



Фоновий сервіс виконує 10 поколінь по 5 кандидатів, здійснюючи звернення до QPU для оцінки кожного кандидата. Після завершення навчені параметри \(\Theta \in \mathbb{R}^{24}\) (по 8 параметрів на кожен із 3 шарів) зберігаються як активна модель для подальшого використання при інференсі.



\begin{center}\rule{0.5\linewidth}{0.5pt}\end{center}



\subsection{8. Ефекти від використання}\label{ux435ux444ux435ux43aux442ux438-ux432ux456ux434-ux432ux438ux43aux43eux440ux438ux441ux442ux430ux43dux43dux44f}



Використання запропонованого рішення забезпечує формування стандартизованого показника CME (Вн), що гарантує відтворюваність та порівнянність результатів оцінювання когнітивного стану між різними сесіями, користувачами та апаратними конфігураціями. Метаевристична оптимізація параметрів квантової схеми та ресурсних режимів \((S, D)\) дозволяє автоматично балансувати між якістю оцінювання та обчислювальною вартістю звернень до QPU, зменшуючи витрати без втрати якості (в допустимій зоні 2\%). Гібридний квантово-класичний інференс із конфігурованим зважуванням квантової та нейромережевої моделей підвищує стійкість оцінювання та забезпечує можливість бенчмаркінгу квантового підходу на реальних даних. Архітектура системи з асинхронним збереженням та автоматичним евристичним маркуванням забезпечує неблокуючу потокову обробку та накопичення розмічених наборів даних для подальшого вдосконалення моделей без ручного маркування.



\begin{center}\rule{0.5\linewidth}{0.5pt}\end{center}



\subsection{9. Літературні посилання}\label{ux43bux456ux442ux435ux440ux430ux442ux443ux440ux43dux456-ux43fux43eux441ux438ux43bux430ux43dux43dux44f}



{[}1{]} Csikszentmihalyi, M. (1990) \emph{Flow: The Psychology of Optimal Experience}. New York: Harper \& Row.



{[}2{]} Pope, A.T., Bogart, E.H. and Bartolome, D.S. (1995) `Biocybernetic system evaluates indices of operator engagement in automated task', \emph{Biological Psychology}, 40(1--2), pp.~187--195.



{[}3{]} Freeman, F.G., Mikulka, P.J., Prinzel, L.J. and Scerbo, M.W. (1999) `Evaluation of an adaptive automation system using three EEG indices with a visual tracking task', \emph{Biological Psychology}, 50(1), pp.~61--76.



{[}4{]} Katahira, K., Yamazaki, Y., Yamaoka, C., Ozaki, H., Nakagawa, S. and Nagata, N. (2018) `EEG correlates of the flow state: A combination of increased frontal theta and moderate frontocentral alpha rhythm in the mental arithmetic task', \emph{Frontiers in Psychology}, 9, 300.



{[}5{]} Pérez-Salinas, A., Cervera-Lierta, A., Gil-Fuster, E. and Latorre, J.I. (2020) `Data re-uploading for a universal quantum classifier', \emph{Quantum}, 4, 226.



{[}6{]} Raufi, B. and Longo, L. (2022) `An evaluation of the EEG alpha-to-theta and theta-to-alpha band ratios as indexes of mental workload', \emph{Frontiers in Neuroinformatics}, 16, 861967.



{[}7{]} Cherep, M., Lim, A.F., Thakor, N. and Jung, T.-P. (2022) `Wearable EEG estimation of flow state during video game play', in \emph{Proceedings of the ACM International Joint Conference on Pervasive and Ubiquitous Computing (UbiComp)}. New York: ACM.



{[}8{]} Hang, L., Chen, X. and Wang, Y. (2024) `Neural correlates of flow using single-channel prefrontal EEG', \emph{Scientific Reports}, 14.



{[}9{]} Ahmed, S., Mühl, C. and Kohlhase, M. (2025) `EEG engagement indices and self-reported flow in cognitive games', \emph{Frontiers in Neuroergonomics}, 6.



{[}10{]} Olvera, C., Montiel Ross, O. and Sepúlveda, R. (2024) `Hybrid quantum machine learning for EEG motor imagery classification', \emph{Neural Computing and Applications}, 36, pp.~5263--5281.



{[}11{]} Hernández-Arango, P., Hernández Mejía, C. and Restrepo-Martínez, A. (2024) `QEEGNet: Quantum Machine Learning for Enhanced Electroencephalography Encoding', arXiv:2407.19214.



{[}12{]} Mohammad, A., Krol, A. and Sarkar, A. (2024) `Meta-optimization of quantum-resource usage for variational tasks', \emph{Scientific Reports}, 14, 10890.



{[}13{]} Wu, X., Zhang, C. and Ding, Y. (2024) `Noise-aware quantum job scheduling with shot- and latency-aware resource management', \emph{ACM Transactions on Quantum Computing}, 5(3).



{[}14{]} Sim, S., Johnson, P.D. and Aspuru-Guzik, A. (2019) `Expressibility and entangling capability of parameterized quantum circuits for hybrid quantum-classical algorithms', \emph{Advanced Quantum Technologies}, 2(12), 1900070.



{[}15{]} Havlíček, V., Córcoles, A.D., Temme, K., Harrow, A.W., Kandala, A., Chow, J.M. and Gambetta, J.M. (2019) `Supervised learning with quantum-enhanced feature spaces', \emph{Nature}, 567, pp.~209--212.



{[}16{]} US7865235B2 (2011) \emph{Methods and system for detecting mental states}. United States Patent.



{[}17{]} US20160235324A1 (2016) \emph{Self-calibrating EEG neurofeedback system}. United States Patent Application.



{[}18{]} US20180279899A1 (2018) \emph{Achieving flow state using biofeedback}. United States Patent Application.



{[}19{]} US20210127981A1 (2021) \emph{Determination and correlation of flow states}. United States Patent Application.



{[}20{]} WO2022231846A1 (2022) \emph{In-queue optimizations for quantum processing units}. WIPO Patent Application.



{[}21{]} US11972321B2 (2024) \emph{Co-scheduling quantum computing jobs}. United States Patent.



{[}22{]} US20240045711A1 (2024) \emph{Preempting a quantum program based on shot count}. United States Patent Application.



{[}23{]} BIPM (2019) \emph{The International System of Units (SI)}. 9th edn. Sèvres: Bureau International des Poids et Mesures.



{[}24{]} Jackson, S.A. and Marsh, H.W. (1996) `Development and validation of a scale to measure optimal experience: The Flow State Scale', \emph{Journal of Sport and Exercise Psychology}, 18(1), pp.~17--35.



{[}25{]} Rheinberg, F., Vollmeyer, R. and Engeser, S. (2003) `Die Erfassung des Flow-Erlebens (The assessment of flow experience)', in Stiensmeier-Pelster, J. and Rheinberg, F. (eds.) \emph{Diagnostik von Motivation und Selbstkonzept}. Göttingen: Hogrefe, pp.~261--279.


\end{document}
